\section{Related work and scope}
Communication topology optimization already includes spatial-temporal
pruning \citep{agentprune}, round-specific agent and edge elimination
\citep{agentdropout}, task-conditioned graph generation \citep{gdesigner},
and joint node/edge pruning \citep{agp}. We do not claim that graph editing,
node omission, or learned topology is new. CATS uses attribution to guide
team search \citep{cats}; this differs from our specific measurement of
existing edge credits after an exogenous graph change, but that distinction
alone does not demonstrate novelty or efficacy.

Active task selection for communication-structure optimization already
accounts for selection and training expenditures \citep{activegraph}.
Its version-three appendix also describes overlapping-subset prediction reuse
and explains that updated graph parameters can change later task utility;
the reported main experiments use single-round selection. Our study therefore
does not claim first use of caching, active graph learning, or changing
information value. DIVA directly values information by its effect on a
decision and stops when checking is not worthwhile \citep{diva}. The relevant
comparison here is a common-predictor, complete-cost revalidation experiment,
not a new general theory of value of information.

Information propagation can help or spread errors \citep{propagation}.
Bounded relays trade context reduction against information loss
\citep{bottleneck}, while information-asymmetric collaboration constitutes a
different deployment regime \citep{iagents}. The centralized Bayes dominance
result of \citet{reliability} concerns an ideal decision maker with common
evidence; it is not a theorem that a finite single LLM dominates every
multi-agent implementation. Accordingly, raw-evidence access and inexpensive
single-model and ensemble baselines are explicit controls in this study.
These cited benchmark results have not been independently reproduced here.

\paragraph{Decision value of feedback.}
Choosing observations by the expected improvement in the best posterior
terminal decision is established knowledge-gradient methodology
\citep{frazier2008knowledge}. That work studies independent normal
populations with known common variance. Our ternary paired outcomes and
graph-dependent query costs do not satisfy that setting, and we do not
inherit its optimality guarantees. The one-step information-value criterion
is an implementation foundation, not a novelty claim.
